\documentclass[11pt]{report}

\usepackage[margin=1in]{geometry}
\usepackage[T1]{fontenc}
\usepackage{lmodern}
\usepackage{hyperref}
\usepackage{listings}
\usepackage{xcolor}
\usepackage{longtable}
\usepackage{booktabs}

\hypersetup{
  colorlinks=true,
  linkcolor=blue,
  urlcolor=blue,
  citecolor=blue
}

\lstdefinestyle{openhac}{
  basicstyle=\ttfamily\small,
  breaklines=true,
  columns=fullflexible,
  frame=single,
  framerule=0.2pt,
  rulecolor=\color{black!30},
  backgroundcolor=\color{black!2},
}

\title{OpenHaC: Hardware-as-Code\\A Detailed Technical Report}
\author{OpenHaC Project}
\date{\today}

\begin{document}
\maketitle
\tableofcontents

\chapter{Executive summary}

OpenHaC is a Python-first toolchain for defining hardware in code and producing KiCad-compatible outputs.
This report documents:
\begin{itemize}
  \item the architecture and compile pipeline,
  \item artifact outputs and handoff formats,
  \item the quality gates used to move toward fabrication-grade results,
  \item verification hooks (SPICE and SI/PI intent),
  \item reproducibility and release evidence.
\end{itemize}

\chapter{Scope and positioning}

\section{What OpenHaC is}
\begin{itemize}
  \item A compiler from a declarative Python design graph (SKiDL-based) into KiCad project artifacts.
  \item A policy engine: ERC/DRC gates, strictness modes, fabrication profiles, and evidence generation.
\end{itemize}

\section{What OpenHaC is not (yet)}
Some engineering domains cannot be fully proven without measurements and external tools (e.g. EMC compliance).
OpenHaC therefore uses a layered approach: encode intent, enforce hard checks, and export evidence.

\chapter{Architecture overview}

\section{High-level structure}
\begin{itemize}
  \item \textbf{Core:} \texttt{openhac/core/} (\texttt{Board}, \texttt{Module}, \texttt{Component})
  \item \textbf{Compiler:} \texttt{openhac/compiler/} (pipeline phases and generators)
  \item \textbf{Database:} \texttt{openhac/database/} (component catalog, offers, alternates)
\end{itemize}

\section{Core primitives}
\subsection{Board}
\texttt{Board} is the top-level orchestrator. It owns:
\begin{itemize}
  \item geometry: \texttt{size\_mm}, \texttt{layers},
  \item modules and constraints,
  \item compile policy: \texttt{compile\_goal}, \texttt{board\_class}, \texttt{quality\_gates},
  \item optional fabrication intent: keepouts, pours, net-ties, mounting holes.
\end{itemize}

\subsection{Module and Interface}
\texttt{Module} groups parts and exposes named \texttt{Interface}s.
Interfaces are used as the natural boundary for hierarchical schematics and connectivity linting.

\chapter{Compile pipeline}

\section{Pipeline phases}
OpenHaC compiles via ordered phases (ERC/DRC, netlist/BOM, layout, routing, schematic, manifest, release zip).
The pipeline supports a compile loop scaffold that can retry after applying best-effort repairs.

\section{Compile goal: handoff vs fabrication}
\begin{itemize}
  \item \textbf{handoff:} preserve artifacts for review; allow best-effort routing fallback.
  \item \textbf{fabrication:} fail-closed gates (PCB fit, routing thresholds, KiCad PCB DRC).
\end{itemize}

\section{Board classes and profiles}
Board class is a profile input that sets defaults for quality gates and (over time) will drive
placement and routing strategies.

\chapter{Schematic generation}

\section{Symbol resolution}
OpenHaC supports KiCad \texttt{.kicad\_sym} pin coordinate lookup for accurate wire endpoints.
For SKiDL-native parts, OpenHaC can generate a project-local symbol library so KiCad does not render \texttt{?}.

\section{Hierarchy export}
Multi-sheet export can emit:
\begin{itemize}
  \item a root sheet with sheet symbols,
  \item sheet pins derived from \texttt{Module.declare\_interface()} nets,
  \item child sheets with matching hierarchical labels,
  \item per-sheet wiring for nets contained within a module.
\end{itemize}

\section{Schematic lint}
Schematic lint is a best-effort quality gate: e.g. interface nets should be named.
In fabrication mode, lint violations are treated as failures; in handoff mode they are warnings.

\chapter{PCB layout and routing}

\section{Placement}
OpenHaC uses module-level constraints and (optionally) Z3 solving for placement.
Post-processing via \texttt{pcbnew} can emit keepouts, net-ties, mounting holes, and copper pours.

\section{Routing}
Two routing paths exist:
\begin{itemize}
  \item FreeRouting DSN/SES flow (preferred for real routing assistance).
  \item A minimal \texttt{pcbnew}-based fallback that draws straight-line tracks (handoff only).
\end{itemize}

\section{Routing metrics and gates}
Fabrication mode can enforce minimal routing-quality thresholds (e.g. minimum track count) and then run KiCad PCB DRC
(\texttt{kicad-cli pcb drc}) as the authoritative physical-rule gate.

\chapter{Verification}

\section{SPICE}
OpenHaC generates an ngspice-oriented \texttt{.cir} netlist from the same graph.
Optional gates can require that model-required parts have \texttt{Spice\_Subckt} annotations and can run ngspice headlessly.

\section{SI/PI handoff}
OpenHaC records intent for:
\begin{itemize}
  \item differential pairs,
  \item length-match groups,
  \item net roles and merge hints,
  \item stackup references.
\end{itemize}
These are packaged as a consolidated handoff JSON for downstream tools and review.

\chapter{Reproducibility and release evidence}

\section{Determinism}
Deterministic mode aims for byte-stable artifacts where possible (sorted JSON, stable timestamps, stable zip entries).

\section{Evidence bundle}
Evidence artifacts include:
\begin{itemize}
  \item manifest output inventory,
  \item evidence index markdown,
  \item optional attestation metadata.
\end{itemize}

\chapter{Appendix: commands}

\section{Compile}
\begin{lstlisting}[style=openhac]
openhac compile design.py --name out --compile-goal handoff -o build/
openhac compile design.py --name out --compile-goal fabrication -o build/
\end{lstlisting}

\section{Simulate}
\begin{lstlisting}[style=openhac]
openhac simulate design.py --name sim -o build/ --spice-preset op
openhac simulate design.py --name sim -o build/ --run-ngspice
\end{lstlisting}

\end{document}

